\documentclass[12pt]{article}
\usepackage{mathtools}
\usepackage{amssymb}
\usepackage{fontspec}
\setmainfont{EB Garamond}

% Set page size and margins
\usepackage[a4paper,top=2cm,bottom=1.5cm,left=1cm,right=1cm]{geometry}
\usepackage{fancyhdr}
\pagestyle{fancy}
\setlength{\parindent}{0pt}

% Useful packages
\usepackage{minted}
\usepackage{graphicx}
\graphicspath{{./media/}}
\usepackage{subfig}
\usepackage{float}
\usepackage[hidelinks]{hyperref}
\usepackage{siunitx}
\usepackage{tabularx}  
\usepackage{array}     

% Change "Table" to "Πίνακας"
\renewcommand{\tablename}{Πίνακας}

%Change "Figure" to "Ιστόγραμμα"
\renewcommand{\figurename}{Ιστόγραμμα}

\title{\vspace{-2cm} Θεμελιώσεις Κρυπτογραφίας \\
        Project 2}
\author{Μαρινόπουλος Χριστόφορος 10522}
\date{Εαρινό εξάμηνο 2025}

\lhead{Project 2}
\chead{Θεμελιώσεις Κρυπτογραφίας}
\rhead{Μαρινόπουλος Χριστόφορος}


\begin{document}
\maketitle
\renewcommand{\contentsname}{Περιεχόμενα}
\tableofcontents
\newpage
\section*{Άσκηση 1}
\addcontentsline{toc}{section}{Άσκηση 1}

Υλοποιήστε σε όποια γλώσσα προγραμματισμού θέλετε τον αλγόριθμο:  Γρήγορη ύψωση σε δύναμη mod m και κατόπιν υπολογίστε το γινόμενο: \( 2^{1000} \cdot 3^{101} \cdot 5^{47} (\mod{2^{107} -1}) \) κάνοντας χρήση του αλγόριθμου του Karatsuba.
\subsection*{Απάντηση}
\addcontentsline{toc}{subsection}{Απάντηση}

Ο αλγόριθμος γρήγορη ύψωση σε δύναμη mod m υλοποιήθηκε σε \texttt{Python} και είναι η συνάρτηση \texttt{fastmod()} που βρίσκεται στο αρχείο \texttt{Exercise\_1.py}. Επιπλέον, στο ίδιο αρχείο βρίσκεται και η συνάρτηση \texttt{karatsuba()} η οποία υλοποιεί τον αλγόριθμο του Karatsuba.
\newline
\newline
Με την χρήση τους, το τελικό αποτέλεσμα είναι: \textbf{3265040604348432725872932588791}.
\hfill \blacksquare

\section*{Άσκηση 2}
\addcontentsline{toc}{section}{Άσκηση 2}
Ν.δ.ο. το \( 2^m (m \ge 2) \) δεν είναι άθροισμα διαδοχικών θετικών ακεραίων.


\subsection*{Απάντηση}
\addcontentsline{toc}{subsection}{Απάντηση}
Άθροισμα n διαδοχικών θετικών ακεραίων:
\[
a + (a+1) + (a+2) + ... + (a+n-1) = na + \sum_{i=1}^{n-1}i
\]

Έστω ότι το \( 2^m (m \ge 2) \) είναι άθροισμα διαδοχικών θετικών ακεραίων. Άρα:

\[
2^m = \sum_{i=1}^{n-1} i + na, n > 2, a \in \mathbb{N}
\]
\[
2^m = \frac{(n-1)n}{2} + na = n(\frac{n-1}{2} + a)
\]

Σύμφωνα με το θεμελιώδες θεώρημα της αριθμητικής, κάθε μη μηδενικός φυσικός αριθμός αναλύεται σε γινόμενο πρώτων παραγόντων με μοναδικό τρόπο. Επομένως:

\[
n = 2^k, k \in \mathbb{N}
\]

Άρα:

\[
2^{m-k} = \frac{2^k-1}{2} +a
\]

\[
\frac{2^k-1}{2} = 2^{k-1} - \frac{1}{2}
\]

\[
a, 2^{k-1} \in \mathbb{N} , -\frac{1}{2} \notin \mathbb{N}
\]
Δηλαδή:
\[
\frac{2^k-1}{2} +a \notin \mathbb{N}
\]
Άτοπο αφού:
\[
2^{m-k} \in \mathbb{N}
\]

Άρα το \( 2^m (m \ge 2) \) δεν είναι άθροισμα διαδοχικών θετικών ακεραίων
\hfill \blacksquare

\section*{Άσκηση 3}
\addcontentsline{toc}{section}{Άσκηση 3}

Έστω p,q, διαδοχικοί πρώτοι αριθμοί τέτοι, ώστε \( 3 < p < q\). Τότε, κάθε ακέραιος \( r < q \) διαιρεί το \( p!\).

\subsection*{Απάντηση}
\addcontentsline{toc}{subsection}{Απάντηση}
Για \( r \le p \), προφανώς το r διαιρεί το p!. 
\newline
\newline
Για \( r > p \) , το r είναι σύνθετος αφού βρίσκεται ανάμεσα σε δύο διαδοχικούς πρώτους. Επομένως μπορεί να γραφτεί ως ab και έστω \( 2 \le a \le b <r \).
\newline
\newline
Έστω \( p < b\) τότε \( 2p < ab < q\). Άτοπο λόγω της εικασίας του Bertrand. Επομένως b < p.
\newline
\newline
Έστω \( a_i, i \in \{2,..,m-1\} \) οι μη τετριμμένοι διαιρέτες του r όπου \( a_1 = 1, a_m = r \). 
\newline
\newline
Για \( m > 3 \) τότε \( r = a_2 a_{m-1} \) όπου \( a_2<a_{m-1}<p\) επομένως ο r διαιρεί τον p.

Για \( m = 3 \)
τότε \( r = d^2\). Επειδή \( 3 < p\) τότε \(r > 5 > 4 = 2^2\) άρα \( d \ge 3 \). 
\newline
\newline
Για \(d = 3\) τότε \( r = 9, p = 7, q = 11\). \( p! = 5040\) και \( 5040 \mod 9 = 0\). 

Για \( d > 3\) έχουμε \(4d \le d^2 \le 2p\) δηλαδή \(  d< 2d < p \) άρα το r διαιρεί το p!.
\newline
\newline
Τελικώς, αποδείχθηκε πως κάθε ακέραιος \( r < q \) διαιρεί το \( p!\) όταν \( 3 < p < q\) και p,q, διαδοχικοί πρώτοι.
\hfill \blacksquare


\section*{Άσκηση 4}
\addcontentsline{toc}{section}{Άσκηση 4}

Έστω a,b θετικοί ακέραιοι. Αν \( \text{gcd}(a,b) = 1 \), τότε ν.δ.ο:
\begin{enumerate}
    \item Για κάθε ακέραιο c ισχύει: \( \text{gcd}(ac,b) = \text{gcd}(c,b) \). 
    \item \( \text{gcd}(a+b,a-b) \in \{1,2\}\). Ειδικότερα, ν.α.ο αν a,b περιττοί θετικοί ακέραιοι, τότε \( \text{gcd}(a+b,a-b) = 2\).
    \item \( \text{gcd}(M_p,M_q) = 1\), όπου \(M_p = 2^p - 1, M_q = 2^q - 1\), Mersenne ακέραιοι (p,q) πρώτοι με \( p \ne q \).
\end{enumerate}

\subsection*{Απάντηση}
\addcontentsline{toc}{subsection}{Απάντηση}
\begin{enumerate}
    \item  Έστω \( d_1 = \text{gcd}(ac,b) \) και \(d_2 = \text{gcd}(c,b) \). Προφανώς \( d_1 \ge d_2 \).
    \newline
    \newline
    Εάν \( d_1 > d_2 \) τότε το \( d_1 \) διαιρεί το a και \( \text{gcd}(a,b) \neq 1 \), Άτοπο . Επομένως \( d_1 = d_2 \).     
    \item Έστω \( d = \text{gcd}(a+b,a-b) \)
    \newline
    \newline
    Τότε το d διαιρεί το \( (a+b) + (a-b) = 2a \) και το \( (a+b) - (a-b) = 2b \). Επίσης \( \text{gcd}(2a,2b) = 2*\text{gcd}(a,b) = 2 \). Άρα d διαιρεί το 2, επομένως \( d \in \{1,2\} \). 
    \newline
    \newline
    Αν a,b περιττοί τότε a+b και a-b ζυγοί άρα d = 2.
    \item p,q πρώτ́οι άρα: \( \text{gcd}(p,q) = 1 \)
    \newline
    \newline
    \( \text{gcd}(2^p-1,2^q-1) = 2^{\text{gcd}(p,q)}-1 = 2^1 -1 = 1 \).

\end{enumerate}
\hfill \blacksquare

\section*{Άσκηση 5}
\addcontentsline{toc}{section}{Άσκηση 5}

Επαληθεύστε πειραματικά ότι για όλους τους περιττούς ακέραιους \( < 2^{20} \) ισχύει:
\[ \frac{\text{σ}(\text{n})}{n} < \frac{e^\text{γ}}{2}\ln{\ln{n} + \frac{0.74}{\ln{\ln{n}}}} 
\]

\subsection*{Απάντηση}
\addcontentsline{toc}{subsection}{Απάντηση}
Για να επαληθευτεί πειραματικά το παραπάνω αρκεί να υπολογιστεί η συνάρτηση σ(n) για κάθε περιττό \( n < 2^{20} \) όπου σ(i) είναι το άθροισμα των (θετικών) διαιρετών του i. 
\newline
\newline
Ο υπολογισμός του σ(n) γίνεται μέσω της συνάρτησης \texttt{sieve()} η οποία βρίσκεται στο αρχείο \texttt{Exercise\_5.py}. Ο έλεγχος της ανισότητας για κάθε περιττό < \( 2^{20} \)γίνεται μέσω της συνάρτηση \texttt{ineq\_check()} η οποία βρίσκεται στο ίδιο αρχείο. Καλώντας την, βλέπουμε πως \textbf{κάθε περιττός \( 2^{20} \) ικανοποιεί την ανισότητα}.
\hfill \blacksquare
\section*{Άσκηση 6}
\addcontentsline{toc}{section}{Άσκηση 6}

Να βρείτε (με χρήση κατάλληλου κώδικα) τον επόμενο τέλειο αριθμό που είναι \( >8128 \). Πόσο χρόνο χρειάστηκε ο κώδικάς σας για το αποτέλεσμα; 

\subsection*{Απάντηση}
\addcontentsline{toc}{subsection}{Απάντηση}
Ο Ευκλείδης απέδειξε πως οι τέλειοι, άρτιοι, αριθμοί είναι της μορφής \(2^{n-1} (2^n -1) \)
όταν το \( 2^n - 1 \) είναι πρώτος. Οι περιττοί τέλειοι αριθμοί ικάζεται πως είναι της μορφής \( 12m + 1 \) ή \( 36m + 9 \) αλλά μέχρι σήμερα δεν έχει βρεθεί κάποιος. 
\newline
\newline
Λαμβάνοντας υπόψη τα παραπάνω αναπτύχθηκε η συνάρτηση \texttt{next\_perfect(num)} που βρίσκεται στο αρχείο 

\texttt{Exercise\_6.py}, δέχεται ως όρισμα έναν ακέραιο που είναι μικρότερος από τον τέλειος αριθμό που ψάχνουμε και 

λειτουργεί ως εξής :
\begin{enumerate}
    \item Αρχικοποιεί το n στο 2.
    \item Υπολογίζει το \(2^n -1\)
    \item Εάν είναι πρώτος υπολογίζει το \(2^{n-1}(2^n-1) \) και προχωράει στο επόμενο βήμα, εάν δεν είναι, αυξάνει το n κατά 1 και επαναλαμβάνει το προηγούμενο βήμα. 
    \item Εάν το \(2^{n-1}(2^n-1) \) είναι μεγαλύτερο από το όρισμα της συνάρτησης τον επιστρέφει, εάν δεν είναι αυξάνει το n κατά 1 μέχρι να βρεθεί μεγαλύτερος τέλειος αριθμός.
\end{enumerate}

Εκτελώντας τον αλγόριθμο, ο επόμενος τέλειος αριθμός μετά τον 8128 είναι ο \textbf{33550336} και ο χρόνος που χρειάστηκε ο κώδικας για να παράξει το αποτέλεσμα είναι αμελητέος.
\hfill \blacksquare
\section*{Άσκηση 7}
\addcontentsline{toc}{section}{Άσκηση 7}

Ικανοποιεί το Τεστ του Fermat ο αριθμός:  \(3714089895285 * 2^{60000} - 1\);

\subsection*{Απάντηση}
\addcontentsline{toc}{subsection}{Απάντηση}

Το τεστ του Fermat υλοποιήθηκε ως συνάρτηση \texttt{Python} στο αρχείο 
\texttt{Exercise\_7.py} και είναι η συνάρτηση \texttt{fermat\_test()}. Εκτελέστηκε για τον αριθμό  \(3714089895285 * 2^{60000} - 1\) με ένα μεγάλο πλήθος διαφορετικών βάσεων και \textbf{δεν απέτυχε καμία φορά το τεστ}.
\hfill \blacksquare

\section*{Άσκηση 8}
\addcontentsline{toc}{section}{Άσκηση 8}

Βρείτε τον μικρότερο ισχυρό ψευδοπρώτο ως προς τη βάση 32.

\subsection*{Απάντηση}
\addcontentsline{toc}{subsection}{Απάντηση}

Για την επίλυση της άσκηση υλοποιήθηκε η συνάρτηση \texttt{find\_strong\_pseudo} στο αρχείο \texttt{Exercise\_8.py} η οποία δέχεται ως όρισμα μία βάση και βρίσκει τον μικρότερο ισχυρό ψευδοπρώτο της ακολουθώντας την δομή του τεστ Miller-Rabin. 
\newline
\newline
Καλώντας την με όρισμα τον αριθμό 32, το αποτέλεσμα είναι ο αριθμός \textbf{25}.

\hfill \blacksquare
\section*{Άσκηση 9}
\addcontentsline{toc}{section}{Άσκηση 9}

Να υλοποιήσετε (σε όποια γλώσσα προγραμματισμού επιθυμείτε) τον αλγόριθμο δοκιμαστικής διαίρεσης και να 

παραγοντοποιήσετε τους αριθμούς \(2^{62} - 1, 2^{102} - 1\).

\subsection*{Απάντηση}
\addcontentsline{toc}{subsection}{Απάντηση}

O αλγόριθμος δοκιμαστικής διαίρεσης υλοποιήθηκε σε \texttt{Python}, ονομάζεται \texttt{trial\_division(n)}, ακολουθεί την δομή του αλγορίθμου 10.1.1  και βρίσκεται στο αρχείο 
\texttt{Exercise\_9.py}.  
\newline
\newline
Με την εκτέλεση του κώδικα, η παραγοντοποίηση του \( 2^{62} - 1 \) είναι : \textbf{[3, 715827883, 2147483647]} 

και του \( 2^{102} - 1 \) είναι : \textbf{[3, 3, 7, 103, 307, 2143, 2857, 6529, 11119, 43691, 131071]}.


\hfill \blacksquare
\section*{Άσκηση 10}
\addcontentsline{toc}{section}{Άσκηση 10}

Αν \( p = 3571, y = 2763, g = 2\), υπολογίστε τον διακριτό λογάριθμο \( dlog_g(y) = x \) με τον αλγόριθμο ρ του Pollard.

\subsection*{Απάντηση}
\addcontentsline{toc}{subsection}{Απάντηση}

Για να υπολογισθεί ο διακριτός λογάριθμος \( dlog_g(y) = x \)
με τον αλγόριθμο ρ του Pollard υλοποιήθηκε η συνάρτηση 

\texttt{pollard\_d\_log()}, βρίσκεται στο αρχείο \texttt{Exercise\_10.py} και λειτουργεί ως εξής:
\begin{enumerate}
    \item Αρχικοποιεί με τυχαίους ακεραίους τα a, b.
    \item Υπολογίζει \( x = g^ah^b \)
    \item Ορίζει την εσωτερική συνάρτηση μετάβασης
    \[
    f\colon (x,a,b)\;\mapsto\;(x',a',b'),
    \]
    ώστε αν \(x \equiv g^a\,h^b \pmod{p}\), τότε:
    \[
    \begin{cases}
        \text{αν }x \bmod 3 = 0: 
            & x' = x\cdot g \pmod{p},\; a' = a + 1 \pmod{N},\; b' = b,\\
        \text{αν }x \bmod 3 = 1: 
            & x' = x^2 \pmod{p},\; a' = 2a \pmod{N},\; b' = 2b \pmod{N},\\
        \text{αν }x \bmod 3 = 2: 
            & x' = x\cdot h \pmod{p},\; a' = a,\; b' = b + 1 \pmod{N}. 
    \end{cases}
    \]
    \item Αρχικοποιεί slow = f(x, a, b), fast = f(f(x, a, b))
    και τα υπολογίζει εως ότου slow.x = fast.x
    \item Αν gcd(slow.b - fast.b, N) \neq 1 επαναλαμβ́ανετε από το βήμα 1 αλλιώς επιστρέφει τον ζητούμενο αριθμό.
\end{enumerate}

Καλώντας την, ο ζητούμενος διακριτός λογάριθμος είναι \textbf{900}.

\hfill \blacksquare

\section*{Άσκηση 11}
\addcontentsline{toc}{section}{Άσκηση 11}

Θεωρούμε έναν φυσικό αριθμό N και ζητάμε έναν διαιρέτη του. Έστω \( F(X) = (X^2 +1)\mod{N} \) και ας είναι \( X_0 \) μια τυχαία τιμή από το \{2,3, ..., N-1\}. Η γραμμή 6 στον αλγόριθμο: κυκλικός αλγόριθμος του Floyd, τώρα θα είναι: 

\begin{center}
    If \( 1 < \text{gcd}(|X-Y|,N) < N \).
\end{center}
Αν δεν βρεθεί κάποιος μη τετριμμένος διαιρέτης, διαλέγουμε άλλη τιμή στο \( X_0\).

Υλοποιήστε τον αλγόριθμο και βρείτε έναν διαιρέτη του αριθμού \( N = 2^{257} - 1 \).
\subsection*{Απάντηση}
\addcontentsline{toc}{subsection}{Απάντηση}

Ο αλγόριθμος υλοποιήθηκε στην γλώσσα \texttt{Python} και βρίσκεται στο αρχείο \texttt{Exercise\_11.py}. Ονομάζεται \texttt{modifiedFloyd}, δέχεται δύο ορίσματα: 
\begin{enumerate}
    \item f: λίστα με τους συντελεστές του πολυωνύμου που θα χρησιμοποιήσει ο αλγόριθμος
    \item N: Ακέραιος του οποίου ένας μη τετριμμένος διαιρέτης του θα επιστραφεί (αν υπάρχει).
\end{enumerate}
  και λειτουργεί ως εξής:
\begin{enumerate}
    
    \item Διαλέγεται ένα τυχαίο σημείο εκκίνησης με το οποίο αρχικοποιούνται τα x και y. 
    \item Υπολογίζονται τα f(x), f(f(y))
    \item Ενημερώνεται το x = f(x) και το y = f(f(y))
    \item Υπολογίζεται το gcd(|x-y|,N), εάν είναι μεγαλύτερο του 1 και μικρότερο του Ν επιστρέφεται το αποτέλεσμα, εάν δεν είναι επαναλαμβάνεται ο αλγόριθμος από το βήμα 2.
    \item Εάν gcd(|x-y|,N) = Ν διαλέγεται διαφορετικό τυχαίο σημείο εκκίνησης και επαναλαμβάνεται ο αλγόριθμος από το βήμα 2.
\end{enumerate}

Το f(a) υπολογίζεται με την βοηθητική συνάρτηση \texttt{poly} η οποία δέχεται 2 ορίσματα: 
\begin{enumerate}
    \item f: Λίστα με τους συντελεστές του πολυωνύμου f(a).
    \item a: Ακέραιος.
\end{enumerate} 
και επιστρέφει το f(a).
\newline
\newline
Με την συνάρτηση \texttt{modifiedFloyd}, ένας μη τετριμμένος διαιρέτης του 
\( N = 2^{257} - 1 \)  είναι ο: 
\begin{center}
    \textbf{535006138814359}.
\end{center}
\hfill \blacksquare

\section*{Άσκηση 12}
\addcontentsline{toc}{section}{Άσκηση 12}

Δίνεται το δημόσιο κλειδί (N,e) = (11413,19). Βρείτε το ιδιωτικό κλειδί και κατόπιν αποκρυπτογραφήστε το μήνυμα:

\begin{center}
    C = (3203,909,3143,5255,5343,3203,909,9958,5278,5343,9958,5278,
    
     4674,909,9958,792,909,4132,3143,3203,5343,792,3143,4443).
\end{center}
Υποθέστε ότι τα γράμματα στο αρχικό μήνυμα m, αναπαρίστανται από τις ASCII τιμές του (δουλέψτε block by block το C).
\subsection*{Απάντηση}
\addcontentsline{toc}{subsection}{Απάντηση}
Για να αποκρυπτογραφηθεί το μήνυμα C, αρκεί να ακολουθηθούν τα εξής βήματα.
\begin{enumerate}
    \item Υπολογισμός του φ(Ν)
    \item Yπολογισμός του ιδιωτικού κλειδιού d. ( \( d \equiv e^{-1} \mod \text{φ}(N) \) )
    \item Για κάθε αριθμό i στο C, υπολογισμός του \( m = i^d \mod N \) και μετατροπή του m σε χαρακτήρα με βάση την κωδικοποίηση ASCII
\end{enumerate}
Τα βήματα αυτά επιτυγχάνονται από τον αλγ́οριθμο \texttt{decode()} ο οποίος βρίσκεται στο αρχείο \texttt{Exercise\_12.py}.
\newline
\newline
Με την εκτέλεσή του, το αποκρυπτογραφημένο μήνυμα είναι: \textbf{welcowe to the realworld}.

\hfill \blacksquare

\section*{Άσκηση 13}
\addcontentsline{toc}{section}{Άσκηση 13}

Aς θεωρήσουμε (N, e) = (194749497518847283,50736902528669041) και το
κρυπτογραφημένο κείμενο που δίνεται στο 
\href{https://github.com/drazioti/book_crypto/blob/master/public_key_crypto/7.2}{repository}, έχει προκύψει από το textbook-RSA (block by block) και κατόπιν κωδικοποιήθηκε. Εφαρμόστε την επίθεση του
Wiener, για να βρείτε το κλειδί d.
Υποθέτουμε ότι στο αρχικό κείμενο m κάθε χαρακτήρας έχει αντικατασταθεί από την ASCII τιμή του. Τέλος, βρείτε το αρχικό μήνυμα
m.

\subsection*{Απάντηση}
\addcontentsline{toc}{subsection}{Απάντηση}

Η επίθεση του Wiener υλοποιήθηκε στο αρχείο \texttt{Exercise\_13.py} και πραγματοποιείται από την συνάρτηση \texttt{Wiener()}. Όταν η συνάρτηση \texttt{Wiener()} κληθεί με ορίσματα e,N το αποτέλεσμά της είναι:  \textbf{d = 20881}.  
\newline
\newline
Το δοσμένο κρυπτογραφημένο κείμενο είναι πρωτίστως κωδικοποιημένο σε base 64. Αφού πραγματοποιηθεί η αποκωδικοποίησή του, έστω ότι αυτή είναι c, μπορεί να αποκρυπτογραφηθεί από την συνάρτηση \texttt{decode()} του αρχείου \texttt{Exercise\_13.py}. Όταν αυτή καλείται με ορίσματα (c,e,N) το αποτέλεσμα είναι το εξής:
\begin{center}
    \textbf{Just because you are a character doesn't mean that you have character}
\end{center}

\hfill \blacksquare

\newpage

\section*{Άσκηση 14}
\addcontentsline{toc}{section}{Άσκηση 14}

Περιγράψτε τις βασικές διαφορές των SSL/TLS ver. 1.2 και ver. 1.3 . 


\subsection*{Απάντηση}
\addcontentsline{toc}{subsection}{Απάντηση}

\begin{table}[ht]
  \centering
  \small  
  \begin{tabularx}{\textwidth}{|>{\raggedright\arraybackslash}p{3cm}|%
                              >{\raggedright\arraybackslash}X|%
                              >{\raggedright\arraybackslash}X|}
    \hline
    \textbf{Χαρακτηριστικό} & \textbf{TLS 1.2} & \textbf{TLS 1.3} \\
    \hline
    \textbf{Handshake Round Trips} 
      & ≥ 2 γύροι 
      & 1 γύρος (ακόμα και 0-RTT για resumption) \\
    \hline
    \textbf{Cipher Suites} 
      & Block ciphers (AES-CBC, 3DES, RC4) + AEAD 
      & Μόνο AEAD (AES-GCM, ChaCha20-Poly1305) \\
    \hline
    \textbf{Key Exchange} 
      & Στατική RSA, DHE, ECDHE 
      & Μόνο Ephemeral (DHE/ECDHE) → υποχρεωτικά Forward Secrecy \\
    \hline
    \textbf{MAC/Hash} 
      & HMAC-SHA1, HMAC-SHA256, MD5 (Legacy) 
      & HMAC-SHA256, HMAC-SHA384 με HKDF \\
    \hline
    \textbf{Ασφάλεια} 
      & Πολλές legacy/vulnerable επιλογές 
      & Αφαίρεση ευπαθειών, Downgrade protection, 0-RTT (με προσοχή) \\
    \hline
    \textbf{Session Resumption} 
      & Session IDs, Session Tickets 
      & PSK-only tickets, stateless resumption \\
    \hline
    \textbf{Extensions} 
      & Πλήθος, renegotiation, encrypt\_then\_mac 
      & Περιορισμένα, πιο στοχευμένα, λιγότερες επιπλοκές \\
    \hline
    \textbf{Απόδοση (Latency)} 
      & Μεγαλύτερο latency λόγω 2 RT 
      & Μικρότερο latency, δυνατότητα 0-RTT \\
    \hline
  \end{tabularx}
  \caption{Σύγκριση TLS 1.2 vs TLS 1.3}
  \label{tab:tls_comparison}
\end{table}

\hfill \blacksquare

\section*{Άσκηση 15}
\addcontentsline{toc}{section}{Άσκηση 15}

Στείλτε ένα κρυπτογραφημένο μήνυμα στον συγγραφέα (το δημόσιο κλειδί του έχει αναγνωριστικό 0xEB1185F82713D6DF). Από το δημόσιο κλειδί το αναγνωριστικό προκύπτει (σε bash) αν δώσουμε τον παρακάτω κώδικα:

\texttt{\$gpg --list-packets pk.asc | awk '/keyid:/{ print \$2 }'}
\subsection*{Απάντηση}
\addcontentsline{toc}{subsection}{Απάντηση}

Το μήνυμα το οποίο έλαβα είναι το εξής: \textbf{gI@r1 k@p0!3\$ pr0θ3\$M!3\$ e!n@! @πλ\_w\$ pr0\$χhM@t@ g1@ n@ e!π0θoύn !\$τ0r!3\$ π0υ έπr3π3 n@ gr@φt0úN.}.

\hfill \blacksquare
\end{document}